% math.tex — kkt module
% Mathematical results supporting crates/src/kkt/
%
% Labels defined here:
%   lem:kkt, lem:H-quadratic, lem:well-defined,
%   lem:numerical-transition-feasibility, lem:q-error-bound,
%   rem:near-null-lp-search
%
% Sources: thesis/general-case-algorithm.tex,
%   thesis/general-case-algorithm-proof.tex,
%   thesis/appendix-numerical.tex
%
% Notation: A = dual-vertex matrix, H = action matrix,
%   beta = dwell-time weights, mu = Lagrange multiplier (closure),
%   xi = Lagrange multiplier (normalization).

% \input'd from crates/src/math.tex — do not add \documentclass or \begin{document}

% ── Setup: linear system referenced by the lemmas ──
%
% For a polytope K with dual vertices a_i and a pair (S, sigma),
% define the dual-vertex matrix A and action matrix H:
%   A_{id} = a_{sigma(i),d},
%   H_{ij} = omega_0(a_{sigma(j)}, a_{sigma(i)})  for i > j,
%            0                                      for i = j,
%            omega_0(a_{sigma(i)}, a_{sigma(j)})    for i < j.
%
% The linear system is:
%   [ H    A    1 ] [ beta ]   [ 0 ]
%   [ A^T  0    0 ] [ mu   ] = [ 0 ]
%   [ 1^T  0    0 ] [ xi   ]   [ 1 ]
%
% See thesis/general-case-algorithm.tex, Algorithm [alg:ehz].


\begin{lemma}[Double sum as quadratic form]\label{lem:H-quadratic}
% Jörn: text approved (b7a8bc0) — from \begin{lemma} to \end{proof}
% Notation updated to dual-vertex form; math unchanged
\[
  \sum_{1 \leq j < i \leq |S|}
    \beta_{\sigma(i)}\, \beta_{\sigma(j)}\,
    \omega_0(a_{\sigma(j)},\, a_{\sigma(i)})
  = \tfrac{1}{2}\, \beta^\top H\, \beta.
\]
In particular,
$\nabla_{\beta}
  (\beta^\top H\, \beta) = 2H \beta$.
\end{lemma}

\begin{proof}
For $i > j$,
$H_{ij} = \omega_0(a_{\sigma(j)}, a_{\sigma(i)})$
by definition of~$H$.
Hence the double sum equals
$\sum_{j < i} \beta_{\sigma(i)}\, \beta_{\sigma(j)}\, H_{ij}$.
Since $H$ is symmetric with zero diagonal,
$\sum_{j < i} \beta_{\sigma(i)} \beta_{\sigma(j)} H_{ij}
= \tfrac{1}{2}\, \beta^\top H\, \beta$.
\end{proof}


\begin{lemma}[KKT conditions and the algorithm's linear system]%
\label{lem:kkt}
% Jörn: text approved (b7a8bc0) — from \begin{lemma} to \end{proof}
% Notation updated to symmetric convention: (lambda, nu) -> (mu, xi)
% Notation updated to dual-vertex form; math unchanged
For fixed $(S, \sigma)$, the linear
system is the Lagrange system for
critical points of
$\beta^\top H\, \beta$ on the constraint
manifold
$\{A^\top \beta = 0,\; \mathbf{1}^\top \beta = 1\}$.
\end{lemma}

\begin{proof}
The gradient is
$\nabla_{\beta}
  (\beta^\top H \beta) = 2H \beta$
(Lemma~\ref{lem:H-quadratic}).
The Lagrange conditions
$\nabla(\beta^\top H \beta)
= -A\lambda - \mathbf{1}\nu$
give
$2H\beta + A\lambda + \mathbf{1}\nu = 0$.
Absorbing the factor~$2$ into the multipliers
($\mu = \lambda/2$, $\xi = \nu/2$):
\[
  H \beta + A \mu + \mathbf{1} \xi = 0,
\]
which, together with $A^\top \beta = 0$ and
$\mathbf{1}^\top \beta = 1$, is exactly
the linear system.
\end{proof}


\begin{lemma}[Well-definedness]\label{lem:well-defined}
% Jörn: text approved (06a976c) — statement
% Jörn: text approved (0bbc647) — proof
% Notation updated to symmetric convention: (lambda, nu) -> (mu, xi)
% Notation updated to dual-vertex form; math unchanged
If the linear system has multiple solutions
$(\beta, \mu, \xi)$, they all yield the same value
$\beta^\top H \beta$.
\end{lemma}

\begin{proof}
Write the three block rows of the linear system as
\[
  \text{(I)}\ H\beta + A\mu + \mathbf{1}\xi = 0, \qquad
  \text{(II)}\ A^\top \beta = 0, \qquad
  \text{(III)}\ \mathbf{1}^\top \beta = 1.
\]
For any solution, left-multiply~(I) by~$\beta^\top$:
\[
  \beta^\top H\beta
  + \underbrace{(A^\top\beta)^\top}_{=\,0\text{ by (II)}}\mu
  + \underbrace{(\mathbf{1}^\top\beta)}_{=\,1\text{ by (III)}}\xi
  = 0,
\]
so $\beta^\top H\beta = -\xi$.
Now let $(\delta\beta, \delta\mu, \delta\xi)$ be any element
of the null space, so that
\[
  \text{(I$'$)}\ H\,\delta\beta + A\,\delta\mu + \mathbf{1}\,\delta\xi = 0, \qquad
  \text{(II$'$)}\ A^\top\delta\beta = 0, \qquad
  \text{(III$'$)}\ \mathbf{1}^\top\delta\beta = 0.
\]
Left-multiply~(I) by~$\delta\beta^\top$:
\[
  \delta\beta^\top H\beta
  + \underbrace{(A^\top\delta\beta)^\top}_{=\,0\text{ by (II$'$)}}\mu
  + \underbrace{(\mathbf{1}^\top\delta\beta)}_{=\,0\text{ by (III$'$)}}\xi
  = 0.
\]
Left-multiply~(I$'$) by~$\delta\beta^\top$:
\[
  \delta\beta^\top H\,\delta\beta
  + \underbrace{(A^\top\delta\beta)^\top}_{0}\delta\mu
  + \underbrace{(\mathbf{1}^\top\delta\beta)}_{0}\delta\xi
  = 0.
\]
Therefore
\[
  (\beta + \delta\beta)^\top H(\beta + \delta\beta)
  = \beta^\top H\beta
    + 2\,\delta\beta^\top H\beta
    + \delta\beta^\top H\,\delta\beta
  = -\xi. \qedhere
\]
\end{proof}


\begin{lemma}[Transition feasibility]%
\label{lem:numerical-transition-feasibility}
% Source: thesis/appendix-numerical.tex
Let $K \subset \mathbb{R}^4$ be a convex polytope with
facets $F_1, \ldots, F_m$ and dual vertices~$a_k$.
A \emph{transition} $F_i \to F_j$ is a segment of a
Reeb trajectory that travels with velocity~$R_i$
for positive time,
then immediately with velocity~$R_j$ for positive time.
A transition $F_i \to F_j$ exists if and only if
\begin{enumerate}
\item $\omega_0(a_i, a_j) \ge 0$, and
\item $\exists\, x \in F_i \cap F_j$ with
  $\langle a_k, x \rangle < 1$
  for all $k \notin \{i,j\}$ satisfying
  $\omega_0(a_i, a_k) < 0$ or
  $\omega_0(a_j, a_k) > 0$.
\end{enumerate}
\end{lemma}

\begin{proof}
A transition $F_i \to F_j$ exists if and only if
there is a point $x \in F_i \cap F_j$ and
$\varepsilon > 0$ such that
$x + t\, J_0 a_i \in K$ for $t \in [-\varepsilon, 0]$
and $x + t\, J_0 a_j \in K$ for $t \in [0, \varepsilon]$:
since $R_i = 2\, J_0\, a_i$,
traveling with velocity~$R_i$ is equivalent to
moving in direction~$J_0 a_i$.

\emph{Necessity.}
Let $x \in F_i \cap F_j$ realize the transition.
For the approach ($t < 0$) to satisfy
the~$F_j$ constraint:
$\langle a_j, x + t\, J_0 a_i \rangle
= 1 + t\,\omega_0(a_i, a_j) \leq 1$,
which forces $\omega_0(a_i, a_j) \geq 0$
(since $t < 0$).
For the departure ($t > 0$), the~$F_i$ constraint gives
$\langle a_i, x + t\, J_0 a_j \rangle
= 1 - t\,\omega_0(a_i, a_j) \leq 1$,
again requiring $\omega_0(a_i, a_j) \geq 0$.
This proves~(1).

For~(2), suppose $k \notin \{i,j\}$ satisfies
$\omega_0(a_i, a_k) < 0$.
Then
$\langle a_k, x + t\, J_0 a_i \rangle
= \langle a_k, x \rangle + t\,\omega_0(a_i, a_k)$.
Since $t < 0$ and $\omega_0(a_i, a_k) < 0$,
their product $t\,\omega_0(a_i, a_k) > 0$, so
$\langle a_k, x + t\, J_0 a_i \rangle
> \langle a_k, x \rangle$.
If $\langle a_k, x \rangle = 1$, this
exceeds~$1$, violating $x + t\, J_0 a_i \in K$.
Hence $\langle a_k, x \rangle < 1$.
The case $\omega_0(a_j, a_k) > 0$ is analogous
for the departure.

\emph{Sufficiency.}
Given $x \in F_i \cap F_j$ satisfying~(1) and~(2),
we construct $\varepsilon > 0$ witnessing
the transition.
For $k \notin \{i,j\}$ with
$\omega_0(n_i, n_k) \geq 0$ and
$\omega_0(n_j, n_k) \leq 0$:
$\langle n_k, x + t\, J_0 n_i \rangle
\leq \langle n_k, x \rangle \leq h_k$
for $t \leq 0$,
and
$\langle n_k, x + t\, J_0 n_j \rangle
\leq \langle n_k, x \rangle \leq h_k$
for $t \geq 0$.
For the remaining $k \notin \{i,j\}$
(those with $\omega_0(n_i, n_k) < 0$ or
$\omega_0(n_j, n_k) > 0$),
condition~(2) gives
$\langle n_k, x \rangle < h_k$.
Let $\delta$ be the minimum of
$h_k - \langle n_k, x \rangle > 0$
over these~$k$ (or any $\delta > 0$ if there are none),
and choose
$\varepsilon < \delta / \max_k(
|\omega_0(n_i, n_k)| + |\omega_0(n_j, n_k)|)$.
Then all halfspace constraints hold for the
approach and departure, so $x$ together
with~$\varepsilon$ witnesses the transition
$F_i \to F_j$.
\end{proof}


\begin{lemma}[$Q$ error bound from KKT residual]%
\label{lem:q-error-bound}
% Source: thesis/appendix-numerical.tex
Let $M$ be the symmetric KKT matrix
with smallest eigenvalue magnitude
$|\lambda_{\min}| = \min_j |\lambda_j| > 0$,
and let $x_0 = (\beta_0, \mu_0, \xi_0)$ be the exact
solution.
Let $\hat x = (\hat\beta, \hat\mu, \hat\xi)$ be the
numerical solution with residual
\[
  r \;=\; M\hat x - b
  \;=\; \begin{pmatrix}
    H\hat\beta + N\hat\mu + \eta\hat\xi \\
    N^\top\hat\beta \\
    \eta^\top\hat\beta - 1
  \end{pmatrix}
  \;=:\; \begin{pmatrix}r_1 \\ r_2 \\ r_3\end{pmatrix},
\]
where $r_1 \in \mathbb{R}^m$ is the stationarity residual,
$r_2 \in \mathbb{R}^4$ is the constraint residual, and
$r_3 \in \mathbb{R}$ is the normalisation residual.
Define the \emph{residual-corrected objective}
\[
  \tilde Q
  \;=\; Q(\hat\beta)
    + \bigl(r_2^\top\hat\mu + r_3\,\hat\xi\bigr)
\]
and the \emph{error bound}
\[
  E \;=\; \frac{9}{2}\,\frac{\|r\|^2}{|\lambda_{\min}|}.
\]
Then $|Q(\beta_0) - \tilde Q| \leq E$.
\end{lemma}

\begin{proof}
Write $\delta x = \hat x - x_0$ with components
$(\delta\beta, \delta\mu, \delta\xi)$.
Since $M x_0 = b$, we have $M\,\delta x = r$,
hence $\|\delta x\| \leq \|r\| / |\lambda_{\min}|$ and
each component satisfies the same bound:
\[
  \|\delta\beta\| \leq \frac{\|r\|}{|\lambda_{\min}|}, \quad
  \|\delta\mu\| \leq \frac{\|r\|}{|\lambda_{\min}|}, \quad
  |\delta\xi| \leq \frac{\|r\|}{|\lambda_{\min}|}.
\]

\emph{Step~1: expand $Q(\hat\beta) - Q(\beta_0)$.}
Since $Q(\beta) = \tfrac{1}{2}\,\beta^\top H\,\beta$,
\[
  Q(\hat\beta) - Q(\beta_0)
  = \delta\beta^\top H\,\beta_0
    + \tfrac{1}{2}\,\delta\beta^\top H\,\delta\beta.
\]

\emph{Step~2: rewrite the cross term using KKT.}
The stationarity condition (first block of $M x_0 = b$) is
$H\beta_0 + N\mu_0 + \eta\xi_0 = 0$.
Hence
\[
  \delta\beta^\top H\,\beta_0
  = -\delta\beta^\top(N\mu_0 + \eta\xi_0)
  = -(N^\top\delta\beta)^\top \mu_0
    - (\eta^\top\delta\beta)\,\xi_0.
\]
The constraint residual blocks give
$N^\top\delta\beta = N^\top\hat\beta = r_2$
(since $N^\top\beta_0 = 0$ by the exact constraint)
and similarly $\eta^\top\delta\beta = r_3$, hence
\[
  -\delta\beta^\top H\,\beta_0
  = r_2^\top\mu_0 + r_3\,\xi_0.
\]
Substituting into Step~1 and rearranging:
\[
  Q(\beta_0) = Q(\hat\beta) + (r_2^\top\mu_0 + r_3\,\xi_0)
    - \tfrac{1}{2}\,\delta\beta^\top H\,\delta\beta.
\]

\emph{Step~3: substitute numerical multipliers.}
Write $\mu_0 = \hat\mu - \delta\mu$ and
$\xi_0 = \hat\xi - \delta\xi$:
\[
  r_2^\top\mu_0 + r_3\,\xi_0
  = (r_2^\top\hat\mu + r_3\,\hat\xi)
    - (r_2^\top\delta\mu + r_3\,\delta\xi).
\]
The first term is the computable correction.
The second is a second-order remainder:
\[
  |r_2^\top\delta\mu + r_3\,\delta\xi|
  \leq \|r_2\|\,\|\delta\mu\| + |r_3|\,|\delta\xi|
  \leq 2\,\frac{\|r\|^2}{|\lambda_{\min}|},
\]
using $\|r_2\|, |r_3| \leq \|r\|$
and the component bounds.

\emph{Step~4: bound the quadratic term
  using the KKT block structure.}
Expanding $\delta x^\top M\,\delta x$
via the block structure
$M = \bigl[\begin{smallmatrix}
  H & N & \eta \\
  N^\top & 0 & 0 \\
  \eta^\top & 0 & 0
\end{smallmatrix}\bigr]$
gives
\[
  \delta\beta^\top H\,\delta\beta
  = \delta x^\top M\,\delta x
    - 2\,r_2^\top\delta\mu
    - 2\,r_3\,\delta\xi,
\]
where we used $N^\top\delta\beta = r_2$
and $\eta^\top\delta\beta = r_3$ from Step~2.
For the first term,
\[
  |\delta x^\top M\,\delta x|
  = |r^\top M^{-1} r|
  = \Bigl|\sum_j \frac{(v_j^\top r)^2}{\lambda_j}\Bigr|
  \leq \frac{\|r\|^2}{|\lambda_{\min}|}.
\]
Combining with the component bounds
for $\|\delta\mu\|$ and $|\delta\xi|$:
\[
  |\delta\beta^\top H\,\delta\beta|
  \leq \frac{\|r\|^2}{|\lambda_{\min}|}
    + 4\,\frac{\|r\|^2}{|\lambda_{\min}|}
  = \frac{5\,\|r\|^2}{|\lambda_{\min}|}.
\]

\emph{Conclusion.}
\[
  |Q(\beta_0) - \tilde Q|
  \leq 2\,\frac{\|r\|^2}{|\lambda_{\min}|}
    + \frac{5}{2}\,\frac{\|r\|^2}{|\lambda_{\min}|}
  = \frac{9}{2}\,\frac{\|r\|^2}{|\lambda_{\min}|}
  = E.
  \qedhere
\]
\end{proof}


\begin{remark}[Near-null LP search for admissible~$\hat\beta$]%
\label{rem:near-null-lp-search}
% Source: saddle_point_solver.rs, beta_feasibility.rs
% Jörn: math approved — Type C impossibility (Type C bullet only), 2026-03-22
% [TODO: JÖRN - verify Type A filtering rationale and Type B bounded-LP argument]

\noindent\textbf{Setting.}
When the augmented KKT matrix~$M$ is numerically rank-deficient
(eigenvalues below threshold~$\tau$),
the pseudoinverse~$\hat\beta$ may have near-zero or negative components.
The true~$\beta^*$ lies approximately in the affine set
$\hat\beta + \operatorname{span}\{
  (\hat v_{j_1})_\beta, \ldots,
  (\hat v_{j_k})_\beta\}$,
where~$(\hat v_j)_\beta \in \mathbb{R}^m$ denotes the
$\beta$-block of the $j$-th discarded eigenvector.

\medskip\noindent\textbf{LP formulation.}
The implementation searches this affine set via
\[
  \max\; t \quad \text{s.t.}\quad
  \hat\beta_i + \sum_{\ell=1}^{k} \alpha_\ell\,
    (\hat v_{j_\ell})_{\beta,i} \;\geq\; t
  \quad \text{for all } i \in S,
\]
over $\alpha \in \mathbb{R}^k$ and $t \in \mathbb{R}$.
If $t^* > 0$, the reconstructed~$\beta$ has all components positive.

\medskip\noindent\textbf{Eigenvector classification.}
Not all discarded eigenvectors contribute usefully to the LP.
The $(m{+}5)$-dimensional eigenvectors split into blocks
$\hat v_j = ((\hat v_j)_\beta,\; (\hat v_j)_\mu,\; (\hat v_j)_\xi)$.
From $M\hat v_j = \hat\lambda_j \hat v_j$, reading the constraint rows:
$N^\top (\hat v_j)_\beta = \hat\lambda_j\, (\hat v_j)_\mu$
and
$\eta^\top (\hat v_j)_\beta = \hat\lambda_j\, (\hat v_j)_\xi$.
Three types arise:

\begin{itemize}
\item \textbf{Type~A} ($\|(\hat v_j)_\beta\| \approx 0$):
  Content lives in the multiplier blocks $(\mu, \xi)$.
  Including these in the LP creates free variables with
  near-zero coefficients, causing spurious unboundedness.
  \emph{Filtered out} before the LP
  (threshold $\|(\hat v_j)_\beta\| < 10^{-10}$).

\item \textbf{Type~B} ($\|(\hat v_j)_\beta\| = O(1)$,
  constraint violation $O(|\hat\lambda_j|)$):
  Approximately constraint-preserving.
  Since $\eta^\top (\hat v_j)_\beta \approx 0$
  with $\eta_i > 0$, the direction has mixed signs,
  so the LP is bounded.
  Handled correctly by the LP.

\item \textbf{Type~C} ($\|(\hat v_j)_\beta\| = O(1)$,
  constraint violation $O(1)$):
  \emph{Cannot occur.}
  From $M\hat v_j = \hat\lambda_j \hat v_j$,
  the constraint rows give
  $N^\top (\hat v_j)_\beta = \hat\lambda_j\, (\hat v_j)_\mu$
  and
  $\eta^\top (\hat v_j)_\beta = \hat\lambda_j\, (\hat v_j)_\xi$.
  Since $\|\hat v_j\| = 1$,
  the constraint violation per unit~$\alpha$ is
  $\leq |\hat\lambda_j| \leq \tau$.
  For discarded eigenvectors ($|\hat\lambda_j| \leq \tau = 10^{-3}$),
  this is $O(\tau)$, not $O(1)$.
  The implementation includes a runtime assertion that panics
  if a Type~C direction is ever encountered.
\end{itemize}
\end{remark}

